\documentclass[11pt]{article}

\usepackage[a4paper,margin=1in]{geometry}
\usepackage{amsmath,amsthm,amssymb}
\usepackage{stmaryrd}% Provides the Kulkarni--Nomizu product symbol \owedge.
\usepackage[most]{tcolorbox}
\usepackage{xcolor}
\usepackage{hyperref}

% ---------- Hyperref setup ----------
\hypersetup{
  colorlinks=true,
  linkcolor=blue!60!black,
  citecolor=blue!60!black,
  urlcolor=blue!60!black
}

% ---------- Paragraph formatting ----------
\setlength{\parindent}{0pt}
\setlength{\parskip}{0.5em}

% ---------- Theorem environments ----------
\newtheorem{theorem}{Theorem}[section]
\newtheorem{lemma}[theorem]{Lemma}
\newtheorem{proposition}[theorem]{Proposition}
\newtheorem{definition}[theorem]{Definition}
\newtheorem{remark}[theorem]{Remark}

% ---------- Colored boxes ----------

% Theorem (blue)
\tcolorboxenvironment{theorem}{
  enhanced,
  breakable,
  colback=blue!3,
  colframe=blue!70!black,
  boxrule=0.8pt,
  arc=2pt,
  left=6pt,
  right=6pt,
  top=6pt,
  bottom=6pt
}

% Lemma (purple)
\tcolorboxenvironment{lemma}{
  enhanced,
  breakable,
  colback=violet!3,
  colframe=violet!70!black,
  boxrule=0.8pt,
  arc=2pt,
  left=6pt,
  right=6pt,
  top=6pt,
  bottom=6pt
}

% Proposition (blue-gray)
\tcolorboxenvironment{proposition}{
  enhanced,
  breakable,
  colback=cyan!3,
  colframe=cyan!60!black,
  boxrule=0.8pt,
  arc=2pt,
  left=6pt,
  right=6pt,
  top=6pt,
  bottom=6pt
}

% Definition (light blue)
\tcolorboxenvironment{definition}{
  enhanced,
  breakable,
  colback=blue!2,
  colframe=blue!50!black,
  boxrule=0.6pt,
  arc=2pt,
  left=6pt,
  right=6pt,
  top=6pt,
  bottom=6pt
}

% Remark (gray)
\tcolorboxenvironment{remark}{
  enhanced,
  breakable,
  colback=gray!5,
  colframe=gray!70!black,
  boxrule=0.6pt,
  arc=2pt,
  left=6pt,
  right=6pt,
  top=6pt,
  bottom=6pt
}

% ---------- Notation ----------
\newcommand{\Tr}{\operatorname{tr}}
\newcommand{\Rm}{\operatorname{Rm}}
\newcommand{\Ric}{\operatorname{Ric}}
\newcommand{\Sc}{\operatorname{Sc}}
\newcommand{\II}{\mathrm{I\!I}}

% ---------- Document ----------
\begin{document}

\section{Conformal Changes of Riemannian Metrics}

We assume that the reader is familiar with the basic notions of smooth
manifolds, Riemannian manifolds, Riemannian metrics, and Levi--Civita
connections.

\subsection{Introduction}

We study the transformation laws for geometric quantities under a conformal
change of metric. Let \((M^n,g)\) be an \(n\)-dimensional Riemannian manifold,
and consider the conformally related metric
\[
\widetilde g=e^{2f}g,
\]
where \(f\in C^\infty(M)\).

The Levi--Civita connection, Riemann curvature tensor, Ricci curvature
tensor, and scalar curvature associated with \(g\) will be denoted by
\[
\nabla,\qquad
\Rm,\qquad
\Ric,\qquad
\Sc,
\]
respectively. The corresponding quantities with respect to
\(\widetilde g\) will be denoted by
\[
\widetilde{\nabla},\qquad
\widetilde{\Rm},\qquad
\widetilde{\Ric},\qquad
\widetilde{\Sc}.
\]

\subsection{Sign Conventions}

Our convention for the Riemann curvature tensor is
\[
\Rm(X,Y)Z
=
\nabla_Y\nabla_XZ
-
\nabla_X\nabla_YZ
+
\nabla_{[X,Y]}Z.
\]

The Ricci curvature tensor is defined by
\[
\Ric(X,Y)
=
\sum_{i=1}^n
g\bigl(\Rm(X,e_i)Y,e_i\bigr),
\]
where \(\{e_i\}_{i=1}^n\) is a local \(g\)-orthonormal frame.

The scalar curvature is defined by
\[
\Sc
=
\Tr_g\Ric
=
\sum_{i=1}^n \Ric(e_i,e_i).
\]
Equivalently,
\[
\Sc
=
\sum_{i,j=1}^n
g\bigl(\Rm(e_i,e_j)e_i,e_j\bigr).
\]

With these conventions, a Riemannian manifold of constant sectional
curvature \(K\) satisfies
\[
\Ric=(n-1)Kg,
\qquad
\Sc=n(n-1)K.
\]

For a smooth function \(f\), its Hessian is defined by
\[
\nabla^2f(X,Y)
=
g(\nabla_X\nabla f,Y),
\]
and the Laplace--Beltrami operator is defined by
\[
\Delta f
=
\Tr_g(\nabla^2f).
\]

\subsection{Hypersurface Conventions}

Let \(\Sigma^{n-1}\subset M^n\) be an immersed hypersurface. With respect
to \(g\), denote its unit normal, second fundamental form, and mean curvature
by
\[
\nu,\qquad
\II,\qquad
H,
\]
respectively.

We adopt the convention
\[
\II(X,Y)
=
g(\nabla_X\nu,Y),
\qquad
X,Y\in T\Sigma,
\]
and define the mean curvature by
\[
H
=
\Tr_{g_\Sigma}\II,
\]
where \(g_\Sigma\) is the metric induced on \(\Sigma\) by \(g\). Thus, \(H\)
is the trace of the second fundamental form, rather than its normalized
average.

The corresponding quantities computed with respect to \(\widetilde g\)
will be denoted by
\[
\widetilde{\nu},\qquad
\widetilde{\II},\qquad
\widetilde H.
\]

Our goal is to derive the transformation laws for these intrinsic and
extrinsic geometric quantities under the conformal change
\[
g\longmapsto \widetilde g=e^{2f}g.
\]

\subsection{Conformal Transformation Formulas}

\begin{definition}[Kulkarni--Nomizu product]
Let \(P\) and \(Q\) be symmetric \((0,2)\)-tensors on \(M\).
Their Kulkarni--Nomizu product is the \((0,4)\)-tensor
\(P\owedge Q\) defined by
\[
\begin{aligned}
(P\owedge Q)(X,Y,Z,W)
={}&
P(X,Z)Q(Y,W)
+
P(Y,W)Q(X,Z)
\\
&-
P(X,W)Q(Y,Z)
-
P(Y,Z)Q(X,W).
\end{aligned}
\]
\end{definition}

We regard the Riemann curvature tensor as a \((0,4)\)-tensor by setting
\[
\Rm(X,Y,Z,W)
=
g\bigl(\Rm(X,Y)Z,W\bigr).
\]

\begin{proposition}[Conformal transformation formulas]
Let
\[
\widetilde g=e^{2f}g,
\qquad
f\in C^\infty(M).
\]
Then, for any smooth vector fields \(X,Y\) on \(M\), the Levi--Civita
connections of \(g\) and \(\widetilde g\) satisfy
\[
\widetilde{\nabla}_X Y
=
\nabla_XY
+
X(f)Y
+
Y(f)X
-
g(X,Y)\nabla f.
\]

Regarding the Riemann curvature tensors as \((0,4)\)-tensors, one has
\[
\boxed{
\widetilde{\Rm}
=
e^{2f}
\left[
\Rm
-
g\owedge
\left(
\nabla^2f
-
df\otimes df
+
\frac12|\nabla f|_g^2\, g
\right)
\right].
}
\]

Equivalently, for any smooth vector fields \(X,Y,Z,W\) on \(M\),
\[
\begin{aligned}
\widetilde{\Rm}(X,Y,Z,W)
=
e^{2f}\Bigl\{
&\Rm(X,Y,Z,W)
\\
&-g(X,Z)\nabla^2f(Y,W)
-g(Y,W)\nabla^2f(X,Z)
\\
&+g(X,W)\nabla^2f(Y,Z)
+g(Y,Z)\nabla^2f(X,W)
\\
&+g(X,Z)Y(f)W(f)
+g(Y,W)X(f)Z(f)
\\
&-g(X,W)Y(f)Z(f)
-g(Y,Z)X(f)W(f)
\\
&-|\nabla f|_g^2
\bigl(
g(X,Z)g(Y,W)
-
g(X,W)g(Y,Z)
\bigr)
\Bigr\}.
\end{aligned}
\]

The Ricci curvature tensors satisfy
\[
\boxed{
\widetilde{\Ric}
=
\Ric
-
(n-2)
\left(
\nabla^2f
-
df\otimes df
+
\frac12|\nabla f|_g^2\, g
\right)
-
\left(
\Delta f+\frac{n-2}{2}|\nabla f|_g^2
\right)g
.
}
\]

The scalar curvatures satisfy
\[
\boxed{
\widetilde{\Sc}
=
e^{-2f}
\left[
\Sc
-
2(n-1)
\left(
\Delta f
+
\frac{n-2}{2}|\nabla f|_g^2
\right)
\right].
}
\]

Let \(\Sigma^{n-1}\subset M^n\) be an immersed hypersurface. Then
\[
\boxed{
\widetilde{\nu}=e^{-f}\nu.
}
\]

For all \(X,Y\in T\Sigma\), the second fundamental forms satisfy
\[
\boxed{
\widetilde{\II}(X,Y)
=
e^f
\left[
\II(X,Y)+\nu(f)g(X,Y)
\right].
}
\]

Finally, the mean curvatures satisfy
\[
\boxed{
\widetilde H
=
e^{-f}
\left[
H+(n-1)\nu(f)
\right].
}
\]
\end{proposition}

\end{document}
